% --- Archon physics formalization source begin ---
% archon:physics
% archon:covers PhyXMiniProblems/problem_phyx_mini_0969.lean
% archon:source-report reports/phyx_mini/problem_phyx_mini_0969.source.json

\chapter{Physics problem phyx\_mini\_0969}
\label{ch:PhyXMiniProblems_problem_phyx_mini_0969}

\paragraph{Problem source.}
\#\# Physical scenario

A cylindrical shell with radius \textbackslash{}( R \textbackslash{}) and length \textbackslash{}( W \textbackslash{}) carries a uniform charge \textbackslash{}( Q \textbackslash{}) and rotates about its axis with angular speed \textbackslash{}( \textbackslash{}omega \textbackslash{}). The center of the cylinder lies at the origin \textbackslash{}( O \textbackslash{}) and its axis is coincident with the \textbackslash{}( x \textbackslash{})-axis.

\#\# Question

A cylindrical shell of radius \textbackslash{}( R \textbackslash{}), length \textbackslash{}( W \textbackslash{}), and total charge \textbackslash{}( Q \textbackslash{}) rotates about the \textbackslash{}( x \textbackslash{})-axis with angular speed \textbackslash{}( \textbackslash{}omega \textbackslash{}), and is centered at the origin \textbackslash{}( O \textbackslash{}).

\#\# Answer choices

- A: A: \textbackslash{}frac\{\textbackslash{}mu\_0\}\{2\textbackslash{}pi\} \textbackslash{}cdot \textbackslash{}frac\{Q\textbackslash{}omega\}\{\textbackslash{}sqrt\{W\textasciicircum{}2 + R\textasciicircum{}2\}\} \textbackslash{}hat\{i\}

- B: B: \textbackslash{}frac\{\textbackslash{}mu\_0\}\{2\textbackslash{}pi\} \textbackslash{}cdot \textbackslash{}frac\{Q\textbackslash{}omega\}\{\textbackslash{}sqrt\{W\textasciicircum{}2 + 4R\textasciicircum{}2\}\} \textbackslash{}hat\{i\}

- C: C: \textbackslash{}frac\{\textbackslash{}mu\_0\}\{2\textbackslash{}pi R\} \textbackslash{}cdot Q\textbackslash{}omega\textbackslash{}

- D: D: \textbackslash{}frac\{\textbackslash{}mu\_0 Q \textbackslash{}omega\}\{2\textbackslash{}pi R\} \textbackslash{}hat\{i\}

\#\# Figure caption (auxiliary; use the image as primary evidence)

The image depicts a 3D cylindrical object seen from the side. Here's a detailed description of the image content:

1. **Cylinder**: A semi-transparent, red-colored cylinder is shown. There are annotations on the cylinder:
   - "Q" is labeled on the left side.
   - "R" is labeled as the radial distance to the edge.
   - "W" indicates the length or width of the cylinder.

2. **Axes**:
   - A horizontal line labeled "x" runs through the center of the cylinder.
   - A point labeled "O" is at the intersection of the horizontal axis with the vertical axis of symmetry of the cylinder.

3. **Additional Objects**:
   - A thin, vertical section within the cylinder is highlighted with a purple line, indicating a small element.
   - The width of this section is labeled "dx."
   - A label "dI" is associated with an arrow pointing downwards along this section, suggesting a current or flow direction.

4. **Arrows**:
   - A black arrow labeled "ω" points downward on the side of the cylinder, possibly indicating rotation or angular velocity.

The image appears to illustrate a concept possibly related to electromagnetism or physics, such as the magnetic field in a cylindrical conductor.

\#\# Dataset metadata

Magnetism; Physical Model Grounding Reasoning, Spatial Relation Reasoning

\paragraph{Recorded answer/context.}
B: B: \textbackslash{}frac\{\textbackslash{}mu\_0\}\{2\textbackslash{}pi\} \textbackslash{}cdot \textbackslash{}frac\{Q\textbackslash{}omega\}\{\textbackslash{}sqrt\{W\textasciicircum{}2 + 4R\textasciicircum{}2\}\} \textbackslash{}hat\{i\}

\paragraph{Figure/image path.}
/root/proposal\_for\_physic/science-mango/phyx\_mini\_run/phyx\_data/test\_image/969.png

\paragraph{Formalization target.}
create a compiling Lean file with sorry bodies at `PhyXMiniProblems/problem\_phyx\_mini\_0969.lean`. The Lean declarations must preserve the physical quantities, dimensions or dimensional roles, figure labels, governing-law hypotheses, and final relation expressed by this problem.
Use Mathlib/Physlib names found through LeanExplore where available. If a physics API is missing, introduce faithful local abstractions rather than scalar placeholder aliases.

\begin{theorem}[Physics formalization target]
\label{thm:physics:phyx_mini_0969:target}
\lean{PhyXMiniProblems.ProblemPhyXMini0969.problem_phyx_mini_0969}
\uses{def:physics:phyx-mini-0969:phyxminiproblems-problemphyxmini0969-lengthinmeters, def:physics:phyx-mini-0969:phyxminiproblems-problemphyxmini0969-chargeincoulombs, def:physics:phyx-mini-0969:phyxminiproblems-problemphyxmini0969-angularspeedinradianspersecond, def:physics:phyx-mini-0969:phyxminiproblems-problemphyxmini0969-magneticfluxdensityinteslas, def:physics:phyx-mini-0969:phyxminiproblems-problemphyxmini0969-axisvector, def:physics:phyx-mini-0969:phyxminiproblems-problemphyxmini0969-rotatingchargedcylindricalshellsetup, def:physics:phyx-mini-0969:phyxminiproblems-problemphyxmini0969-matcheswrittenrotatingchargedcylinderscenario, def:physics:phyx-mini-0969:phyxminiproblems-problemphyxmini0969-matchessuppliedrotatingchargedcylinderfigure, def:physics:phyx-mini-0969:phyxminiproblems-problemphyxmini0969-hasphysicalrotatingchargedcylinderparameters, def:physics:phyx-mini-0969:phyxminiproblems-problemphyxmini0969-usescoherentsivacuumpermeability, def:physics:phyx-mini-0969:phyxminiproblems-problemphyxmini0969-satisfiesuniformrotatingsurfacechargemodel, def:physics:phyx-mini-0969:phyxminiproblems-problemphyxmini0969-satisfiesringbiotsavartandsuperposition, def:physics:phyx-mini-0969:phyxminiproblems-problemphyxmini0969-recordeddatasetanswer, def:physics:phyx-mini-0969:phyxminiproblems-problemphyxmini0969-matchesmodeledcenterfield}
The assigned autoformalize agent should translate this physics problem into Lean declarations in the covered file, with theorem and lemma proof bodies written as `by sorry`.
\end{theorem}
\begin{proof}
This is an autoformalization task, not a proof task. Produce statements that can later be proved by the physics prover without weakening the physical model.
\end{proof}

% --- Archon declaration topology begin ---
\section{Declaration topology}
\begin{definition}[angular Speed Dimension]
  \label{def:physics:phyx-mini-0969:phyxminiproblems-problemphyxmini0969-angularspeeddimension}
  \lean{PhyXMiniProblems.ProblemPhyXMini0969.angularSpeedDimension}
  Angular speed has inverse-time dimension; radians are dimensionless
\end{definition}
\begin{proof}
  This is the declared model definition of angular Speed Dimension.
\end{proof}
\begin{definition}[electric Current Dimension]
  \label{def:physics:phyx-mini-0969:phyxminiproblems-problemphyxmini0969-electriccurrentdimension}
  \lean{PhyXMiniProblems.ProblemPhyXMini0969.electricCurrentDimension}
  Electric current has dimension charge per time
\end{definition}
\begin{proof}
  This is the declared model definition of electric Current Dimension.
\end{proof}
\begin{definition}[surface Charge Density Dimension]
  \label{def:physics:phyx-mini-0969:phyxminiproblems-problemphyxmini0969-surfacechargedensitydimension}
  \lean{PhyXMiniProblems.ProblemPhyXMini0969.surfaceChargeDensityDimension}
  Surface charge density has dimension charge per area
\end{definition}
\begin{proof}
  This is the declared model definition of surface Charge Density Dimension.
\end{proof}
\begin{definition}[linear Charge Density Dimension]
  \label{def:physics:phyx-mini-0969:phyxminiproblems-problemphyxmini0969-linearchargedensitydimension}
  \lean{PhyXMiniProblems.ProblemPhyXMini0969.linearChargeDensityDimension}
  Axial linear charge density has dimension charge per length
\end{definition}
\begin{proof}
  This is the declared model definition of linear Charge Density Dimension.
\end{proof}
\begin{definition}[current Per Length Dimension]
  \label{def:physics:phyx-mini-0969:phyxminiproblems-problemphyxmini0969-currentperlengthdimension}
  \lean{PhyXMiniProblems.ProblemPhyXMini0969.currentPerLengthDimension}
  \uses{def:physics:phyx-mini-0969:phyxminiproblems-problemphyxmini0969-electriccurrentdimension}
  Ring current per unit axial length has dimension current per length
\end{definition}
\begin{proof}
  This is the declared model definition of current Per Length Dimension.
\end{proof}
\begin{definition}[vacuum Permeability Dimension]
  \label{def:physics:phyx-mini-0969:phyxminiproblems-problemphyxmini0969-vacuumpermeabilitydimension}
  \lean{PhyXMiniProblems.ProblemPhyXMini0969.vacuumPermeabilityDimension}
  Vacuum permeability has SI dimension N A⁻² = M L C⁻²
\end{definition}
\begin{proof}
  This is the declared model definition of vacuum Permeability Dimension.
\end{proof}
\begin{definition}[magnetic Flux Density Dimension]
  \label{def:physics:phyx-mini-0969:phyxminiproblems-problemphyxmini0969-magneticfluxdensitydimension}
  \lean{PhyXMiniProblems.ProblemPhyXMini0969.magneticFluxDensityDimension}
  Magnetic flux density has the tesla dimension M T⁻¹ C⁻¹
\end{definition}
\begin{proof}
  This is the declared model definition of magnetic Flux Density Dimension.
\end{proof}
\begin{definition}[Length Quantity]
  \label{def:physics:phyx-mini-0969:phyxminiproblems-problemphyxmini0969-lengthquantity}
  \lean{PhyXMiniProblems.ProblemPhyXMini0969.LengthQuantity}
  A nonnegative, unit-independent physical length
\end{definition}
\begin{proof}
  This is the declared model definition of Length Quantity.
\end{proof}
\begin{definition}[Axial Position Quantity]
  \label{def:physics:phyx-mini-0969:phyxminiproblems-problemphyxmini0969-axialpositionquantity}
  \lean{PhyXMiniProblems.ProblemPhyXMini0969.AxialPositionQuantity}
  A signed, unit-independent coordinate along the cylinder axis
\end{definition}
\begin{proof}
  This is the declared model definition of Axial Position Quantity.
\end{proof}
\begin{definition}[Charge Magnitude Quantity]
  \label{def:physics:phyx-mini-0969:phyxminiproblems-problemphyxmini0969-chargemagnitudequantity}
  \lean{PhyXMiniProblems.ProblemPhyXMini0969.ChargeMagnitudeQuantity}
  A nonnegative total electric-charge magnitude
\end{definition}
\begin{proof}
  This is the declared model definition of Charge Magnitude Quantity.
\end{proof}
\begin{definition}[Angular Speed Quantity]
  \label{def:physics:phyx-mini-0969:phyxminiproblems-problemphyxmini0969-angularspeedquantity}
  \lean{PhyXMiniProblems.ProblemPhyXMini0969.AngularSpeedQuantity}
  \uses{def:physics:phyx-mini-0969:phyxminiproblems-problemphyxmini0969-angularspeeddimension}
  A nonnegative angular-speed magnitude
\end{definition}
\begin{proof}
  This is the declared model definition of Angular Speed Quantity.
\end{proof}
\begin{definition}[Electric Current Quantity]
  \label{def:physics:phyx-mini-0969:phyxminiproblems-problemphyxmini0969-electriccurrentquantity}
  \lean{PhyXMiniProblems.ProblemPhyXMini0969.ElectricCurrentQuantity}
  \uses{def:physics:phyx-mini-0969:phyxminiproblems-problemphyxmini0969-electriccurrentdimension}
  A nonnegative conventional-current magnitude
\end{definition}
\begin{proof}
  This is the declared model definition of Electric Current Quantity.
\end{proof}
\begin{definition}[Surface Charge Density Quantity]
  \label{def:physics:phyx-mini-0969:phyxminiproblems-problemphyxmini0969-surfacechargedensityquantity}
  \lean{PhyXMiniProblems.ProblemPhyXMini0969.SurfaceChargeDensityQuantity}
  \uses{def:physics:phyx-mini-0969:phyxminiproblems-problemphyxmini0969-surfacechargedensitydimension}
  A nonnegative surface charge density
\end{definition}
\begin{proof}
  This is the declared model definition of Surface Charge Density Quantity.
\end{proof}
\begin{definition}[Linear Charge Density Quantity]
  \label{def:physics:phyx-mini-0969:phyxminiproblems-problemphyxmini0969-linearchargedensityquantity}
  \lean{PhyXMiniProblems.ProblemPhyXMini0969.LinearChargeDensityQuantity}
  \uses{def:physics:phyx-mini-0969:phyxminiproblems-problemphyxmini0969-linearchargedensitydimension}
  A nonnegative charge per unit axial length
\end{definition}
\begin{proof}
  This is the declared model definition of Linear Charge Density Quantity.
\end{proof}
\begin{definition}[Ring Current Per Length Quantity]
  \label{def:physics:phyx-mini-0969:phyxminiproblems-problemphyxmini0969-ringcurrentperlengthquantity}
  \lean{PhyXMiniProblems.ProblemPhyXMini0969.RingCurrentPerLengthQuantity}
  \uses{def:physics:phyx-mini-0969:phyxminiproblems-problemphyxmini0969-currentperlengthdimension}
  A nonnegative ring current per unit axial length
\end{definition}
\begin{proof}
  This is the declared model definition of Ring Current Per Length Quantity.
\end{proof}
\begin{definition}[Vacuum Permeability Quantity]
  \label{def:physics:phyx-mini-0969:phyxminiproblems-problemphyxmini0969-vacuumpermeabilityquantity}
  \lean{PhyXMiniProblems.ProblemPhyXMini0969.VacuumPermeabilityQuantity}
  \uses{def:physics:phyx-mini-0969:phyxminiproblems-problemphyxmini0969-vacuumpermeabilitydimension}
  A nonnegative magnetic permeability
\end{definition}
\begin{proof}
  This is the declared model definition of Vacuum Permeability Quantity.
\end{proof}
\begin{definition}[Magnetic Flux Density Quantity]
  \label{def:physics:phyx-mini-0969:phyxminiproblems-problemphyxmini0969-magneticfluxdensityquantity}
  \lean{PhyXMiniProblems.ProblemPhyXMini0969.MagneticFluxDensityQuantity}
  \uses{def:physics:phyx-mini-0969:phyxminiproblems-problemphyxmini0969-magneticfluxdensitydimension}
  A nonnegative magnetic-flux-density magnitude
\end{definition}
\begin{proof}
  This is the declared model definition of Magnetic Flux Density Quantity.
\end{proof}
\begin{definition}[Spatial Vector]
  \label{def:physics:phyx-mini-0969:phyxminiproblems-problemphyxmini0969-spatialvector}
  \lean{PhyXMiniProblems.ProblemPhyXMini0969.SpatialVector}
  A coherent-SI three-dimensional vector readout
\end{definition}
\begin{proof}
  This is the declared model definition of Spatial Vector.
\end{proof}
\begin{definition}[coherent SIReadout]
  \label{def:physics:phyx-mini-0969:phyxminiproblems-problemphyxmini0969-coherentsireadout}
  \lean{PhyXMiniProblems.ProblemPhyXMini0969.coherentSIReadout}
  Coherent-SI scalar readout of a nonnegative dimensionful quantity
\end{definition}
\begin{proof}
  This is the declared model definition of coherent SIReadout.
\end{proof}
\begin{definition}[length In Meters]
  \label{def:physics:phyx-mini-0969:phyxminiproblems-problemphyxmini0969-lengthinmeters}
  \lean{PhyXMiniProblems.ProblemPhyXMini0969.lengthInMeters}
  \uses{def:physics:phyx-mini-0969:phyxminiproblems-problemphyxmini0969-lengthquantity, def:physics:phyx-mini-0969:phyxminiproblems-problemphyxmini0969-coherentsireadout}
  Read a physical length in metres
\end{definition}
\begin{proof}
  This is the declared model definition of length In Meters.
\end{proof}
\begin{definition}[axial Position In Meters]
  \label{def:physics:phyx-mini-0969:phyxminiproblems-problemphyxmini0969-axialpositioninmeters}
  \lean{PhyXMiniProblems.ProblemPhyXMini0969.axialPositionInMeters}
  \uses{def:physics:phyx-mini-0969:phyxminiproblems-problemphyxmini0969-axialpositionquantity}
  Read a signed axial coordinate in metres
\end{definition}
\begin{proof}
  This is the declared model definition of axial Position In Meters.
\end{proof}
\begin{definition}[charge In Coulombs]
  \label{def:physics:phyx-mini-0969:phyxminiproblems-problemphyxmini0969-chargeincoulombs}
  \lean{PhyXMiniProblems.ProblemPhyXMini0969.chargeInCoulombs}
  \uses{def:physics:phyx-mini-0969:phyxminiproblems-problemphyxmini0969-chargemagnitudequantity, def:physics:phyx-mini-0969:phyxminiproblems-problemphyxmini0969-coherentsireadout}
  Read a charge magnitude in coulombs
\end{definition}
\begin{proof}
  This is the declared model definition of charge In Coulombs.
\end{proof}
\begin{definition}[angular Speed In Radians Per Second]
  \label{def:physics:phyx-mini-0969:phyxminiproblems-problemphyxmini0969-angularspeedinradianspersecond}
  \lean{PhyXMiniProblems.ProblemPhyXMini0969.angularSpeedInRadiansPerSecond}
  \uses{def:physics:phyx-mini-0969:phyxminiproblems-problemphyxmini0969-angularspeedquantity, def:physics:phyx-mini-0969:phyxminiproblems-problemphyxmini0969-coherentsireadout}
  Read angular speed in radians per second
\end{definition}
\begin{proof}
  This is the declared model definition of angular Speed In Radians Per Second.
\end{proof}
\begin{definition}[current In Amperes]
  \label{def:physics:phyx-mini-0969:phyxminiproblems-problemphyxmini0969-currentinamperes}
  \lean{PhyXMiniProblems.ProblemPhyXMini0969.currentInAmperes}
  \uses{def:physics:phyx-mini-0969:phyxminiproblems-problemphyxmini0969-electriccurrentquantity, def:physics:phyx-mini-0969:phyxminiproblems-problemphyxmini0969-coherentsireadout}
  Read a conventional-current magnitude in amperes
\end{definition}
\begin{proof}
  This is the declared model definition of current In Amperes.
\end{proof}
\begin{definition}[surface Charge Density In Coulombs Per Square Meter]
  \label{def:physics:phyx-mini-0969:phyxminiproblems-problemphyxmini0969-surfacechargedensityincoulombspersquaremeter}
  \lean{PhyXMiniProblems.ProblemPhyXMini0969.surfaceChargeDensityInCoulombsPerSquareMeter}
  \uses{def:physics:phyx-mini-0969:phyxminiproblems-problemphyxmini0969-surfacechargedensityquantity, def:physics:phyx-mini-0969:phyxminiproblems-problemphyxmini0969-coherentsireadout}
  Read surface charge density in coulombs per square metre
\end{definition}
\begin{proof}
  This is the declared model definition of surface Charge Density In Coulombs Per Square Meter.
\end{proof}
\begin{definition}[linear Charge Density In Coulombs Per Meter]
  \label{def:physics:phyx-mini-0969:phyxminiproblems-problemphyxmini0969-linearchargedensityincoulombspermeter}
  \lean{PhyXMiniProblems.ProblemPhyXMini0969.linearChargeDensityInCoulombsPerMeter}
  \uses{def:physics:phyx-mini-0969:phyxminiproblems-problemphyxmini0969-linearchargedensityquantity, def:physics:phyx-mini-0969:phyxminiproblems-problemphyxmini0969-coherentsireadout}
  Read axial charge density in coulombs per metre
\end{definition}
\begin{proof}
  This is the declared model definition of linear Charge Density In Coulombs Per Meter.
\end{proof}
\begin{definition}[ring Current Per Length In Amperes Per Meter]
  \label{def:physics:phyx-mini-0969:phyxminiproblems-problemphyxmini0969-ringcurrentperlengthinamperespermeter}
  \lean{PhyXMiniProblems.ProblemPhyXMini0969.ringCurrentPerLengthInAmperesPerMeter}
  \uses{def:physics:phyx-mini-0969:phyxminiproblems-problemphyxmini0969-ringcurrentperlengthquantity, def:physics:phyx-mini-0969:phyxminiproblems-problemphyxmini0969-coherentsireadout}
  Read ring current per axial metre in amperes per metre
\end{definition}
\begin{proof}
  This is the declared model definition of ring Current Per Length In Amperes Per Meter.
\end{proof}
\begin{definition}[vacuum Permeability In Newtons Per Ampere Squared]
  \label{def:physics:phyx-mini-0969:phyxminiproblems-problemphyxmini0969-vacuumpermeabilityinnewtonsperamperesquared}
  \lean{PhyXMiniProblems.ProblemPhyXMini0969.vacuumPermeabilityInNewtonsPerAmpereSquared}
  \uses{def:physics:phyx-mini-0969:phyxminiproblems-problemphyxmini0969-vacuumpermeabilityquantity, def:physics:phyx-mini-0969:phyxminiproblems-problemphyxmini0969-coherentsireadout}
  Read vacuum permeability in newtons per ampere squared
\end{definition}
\begin{proof}
  This is the declared model definition of vacuum Permeability In Newtons Per Ampere Squared.
\end{proof}
\begin{definition}[magnetic Flux Density In Teslas]
  \label{def:physics:phyx-mini-0969:phyxminiproblems-problemphyxmini0969-magneticfluxdensityinteslas}
  \lean{PhyXMiniProblems.ProblemPhyXMini0969.magneticFluxDensityInTeslas}
  \uses{def:physics:phyx-mini-0969:phyxminiproblems-problemphyxmini0969-magneticfluxdensityquantity, def:physics:phyx-mini-0969:phyxminiproblems-problemphyxmini0969-coherentsireadout}
  Read a magnetic-flux-density magnitude in teslas
\end{definition}
\begin{proof}
  This is the declared model definition of magnetic Flux Density In Teslas.
\end{proof}
\begin{definition}[Coordinate Axis]
  \label{def:physics:phyx-mini-0969:phyxminiproblems-problemphyxmini0969-coordinateaxis}
  \lean{PhyXMiniProblems.ProblemPhyXMini0969.CoordinateAxis}
  The three Cartesian axes, of which x is the cylinder axis
\end{definition}
\begin{proof}
  This is the declared model definition of Coordinate Axis.
\end{proof}
\begin{definition}[axis Vector]
  \label{def:physics:phyx-mini-0969:phyxminiproblems-problemphyxmini0969-axisvector}
  \lean{PhyXMiniProblems.ProblemPhyXMini0969.axisVector}
  \uses{def:physics:phyx-mini-0969:phyxminiproblems-problemphyxmini0969-spatialvector, def:physics:phyx-mini-0969:phyxminiproblems-problemphyxmini0969-coordinateaxis}
  Standard unit vector associated with a Cartesian axis
\end{definition}
\begin{proof}
  This is the declared model definition of axis Vector.
\end{proof}
\begin{definition}[coordinate Origin]
  \label{def:physics:phyx-mini-0969:phyxminiproblems-problemphyxmini0969-coordinateorigin}
  \lean{PhyXMiniProblems.ProblemPhyXMini0969.coordinateOrigin}
  The coordinate origin carrying the label O
\end{definition}
\begin{proof}
  This is the declared model definition of coordinate Origin.
\end{proof}
\begin{definition}[Cylindrical Shell Part]
  \label{def:physics:phyx-mini-0969:phyxminiproblems-problemphyxmini0969-cylindricalshellpart}
  \lean{PhyXMiniProblems.ProblemPhyXMini0969.CylindricalShellPart}
  The physical part of the cylinder on which the charge resides
\end{definition}
\begin{proof}
  This is the declared model definition of Cylindrical Shell Part.
\end{proof}
\begin{definition}[Charge Distribution]
  \label{def:physics:phyx-mini-0969:phyxminiproblems-problemphyxmini0969-chargedistribution}
  \lean{PhyXMiniProblems.ProblemPhyXMini0969.ChargeDistribution}
  Axial distribution of charge on the shell
\end{definition}
\begin{proof}
  This is the declared model definition of Charge Distribution.
\end{proof}
\begin{definition}[Rotation Regime]
  \label{def:physics:phyx-mini-0969:phyxminiproblems-problemphyxmini0969-rotationregime}
  \lean{PhyXMiniProblems.ProblemPhyXMini0969.RotationRegime}
  Time dependence of the rotation
\end{definition}
\begin{proof}
  This is the declared model definition of Rotation Regime.
\end{proof}
\begin{definition}[Rotation Sense]
  \label{def:physics:phyx-mini-0969:phyxminiproblems-problemphyxmini0969-rotationsense}
  \lean{PhyXMiniProblems.ProblemPhyXMini0969.RotationSense}
  Oriented sense of rotation about the displayed x-axis
\end{definition}
\begin{proof}
  This is the declared model definition of Rotation Sense.
\end{proof}
\begin{definition}[Side View Arrow Direction]
  \label{def:physics:phyx-mini-0969:phyxminiproblems-problemphyxmini0969-sideviewarrowdirection}
  \lean{PhyXMiniProblems.ProblemPhyXMini0969.SideViewArrowDirection}
  Vertical direction of an arrow as it appears in the side-view raster
\end{definition}
\begin{proof}
  This is the declared model definition of Side View Arrow Direction.
\end{proof}
\begin{definition}[Figure Label]
  \label{def:physics:phyx-mini-0969:phyxminiproblems-problemphyxmini0969-figurelabel}
  \lean{PhyXMiniProblems.ProblemPhyXMini0969.FigureLabel}
  Literal labels visible in image 969.png
\end{definition}
\begin{proof}
  This is the declared model definition of Figure Label.
\end{proof}
\begin{definition}[Rotating Charged Cylinder Figure]
  \label{def:physics:phyx-mini-0969:phyxminiproblems-problemphyxmini0969-rotatingchargedcylinderfigure}
  \lean{PhyXMiniProblems.ProblemPhyXMini0969.RotatingChargedCylinderFigure}
  \uses{def:physics:phyx-mini-0969:phyxminiproblems-problemphyxmini0969-lengthquantity, def:physics:phyx-mini-0969:phyxminiproblems-problemphyxmini0969-axialpositionquantity, def:physics:phyx-mini-0969:phyxminiproblems-problemphyxmini0969-chargemagnitudequantity, def:physics:phyx-mini-0969:phyxminiproblems-problemphyxmini0969-angularspeedquantity, def:physics:phyx-mini-0969:phyxminiproblems-problemphyxmini0969-electriccurrentquantity, def:physics:phyx-mini-0969:phyxminiproblems-problemphyxmini0969-coordinateaxis, def:physics:phyx-mini-0969:phyxminiproblems-problemphyxmini0969-sideviewarrowdirection, def:physics:phyx-mini-0969:phyxminiproblems-problemphyxmini0969-figurelabel}
  Presentation data read from the primary image. The dimensionful label values are kept separate from the physical setup and calibrated below. The dx and dI fields represent the highlighted narrow circular ring
\end{definition}
\begin{proof}
  This is the declared model definition of Rotating Charged Cylinder Figure.
\end{proof}
\begin{definition}[Rotating Charged Cylindrical Shell Setup]
  \label{def:physics:phyx-mini-0969:phyxminiproblems-problemphyxmini0969-rotatingchargedcylindricalshellsetup}
  \lean{PhyXMiniProblems.ProblemPhyXMini0969.RotatingChargedCylindricalShellSetup}
  \uses{def:physics:phyx-mini-0969:phyxminiproblems-problemphyxmini0969-lengthquantity, def:physics:phyx-mini-0969:phyxminiproblems-problemphyxmini0969-axialpositionquantity, def:physics:phyx-mini-0969:phyxminiproblems-problemphyxmini0969-chargemagnitudequantity, def:physics:phyx-mini-0969:phyxminiproblems-problemphyxmini0969-angularspeedquantity, def:physics:phyx-mini-0969:phyxminiproblems-problemphyxmini0969-electriccurrentquantity, def:physics:phyx-mini-0969:phyxminiproblems-problemphyxmini0969-surfacechargedensityquantity, def:physics:phyx-mini-0969:phyxminiproblems-problemphyxmini0969-linearchargedensityquantity, def:physics:phyx-mini-0969:phyxminiproblems-problemphyxmini0969-ringcurrentperlengthquantity, def:physics:phyx-mini-0969:phyxminiproblems-problemphyxmini0969-vacuumpermeabilityquantity, def:physics:phyx-mini-0969:phyxminiproblems-problemphyxmini0969-magneticfluxdensityquantity, def:physics:phyx-mini-0969:phyxminiproblems-problemphyxmini0969-coordinateaxis, def:physics:phyx-mini-0969:phyxminiproblems-problemphyxmini0969-cylindricalshellpart, def:physics:phyx-mini-0969:phyxminiproblems-problemphyxmini0969-chargedistribution, def:physics:phyx-mini-0969:phyxminiproblems-problemphyxmini0969-rotationregime, def:physics:phyx-mini-0969:phyxminiproblems-problemphyxmini0969-rotationsense, def:physics:phyx-mini-0969:phyxminiproblems-problemphyxmini0969-rotatingchargedcylinderfigure}
  Independent physical data for the rotating shell. The scalar function axialFieldContributionTeslasPerMeter is the coherent-SI contribution density of the ring at coordinate x; its value is constrained by Biot--Savart below
\end{definition}
\begin{proof}
  This is the declared model definition of Rotating Charged Cylindrical Shell Setup.
\end{proof}
\begin{definition}[Matches Written Rotating Charged Cylinder Scenario]
  \label{def:physics:phyx-mini-0969:phyxminiproblems-problemphyxmini0969-matcheswrittenrotatingchargedcylinderscenario}
  \lean{PhyXMiniProblems.ProblemPhyXMini0969.MatchesWrittenRotatingChargedCylinderScenario}
  \uses{def:physics:phyx-mini-0969:phyxminiproblems-problemphyxmini0969-coordinateorigin, def:physics:phyx-mini-0969:phyxminiproblems-problemphyxmini0969-rotatingchargedcylindricalshellsetup}
  Qualitative and geometric information stated in the written problem
\end{definition}
\begin{proof}
  This is the declared model definition of Matches Written Rotating Charged Cylinder Scenario.
\end{proof}
\begin{definition}[Matches Supplied Rotating Charged Cylinder Figure]
  \label{def:physics:phyx-mini-0969:phyxminiproblems-problemphyxmini0969-matchessuppliedrotatingchargedcylinderfigure}
  \lean{PhyXMiniProblems.ProblemPhyXMini0969.MatchesSuppliedRotatingChargedCylinderFigure}
  \uses{def:physics:phyx-mini-0969:phyxminiproblems-problemphyxmini0969-rotatingchargedcylindricalshellsetup}
  Literal image evidence and calibration of all labelled quantities
\end{definition}
\begin{proof}
  This is the declared model definition of Matches Supplied Rotating Charged Cylinder Figure.
\end{proof}
\begin{definition}[Has Physical Rotating Charged Cylinder Parameters]
  \label{def:physics:phyx-mini-0969:phyxminiproblems-problemphyxmini0969-hasphysicalrotatingchargedcylinderparameters}
  \lean{PhyXMiniProblems.ProblemPhyXMini0969.HasPhysicalRotatingChargedCylinderParameters}
  \uses{def:physics:phyx-mini-0969:phyxminiproblems-problemphyxmini0969-lengthinmeters, def:physics:phyx-mini-0969:phyxminiproblems-problemphyxmini0969-axialpositioninmeters, def:physics:phyx-mini-0969:phyxminiproblems-problemphyxmini0969-chargeincoulombs, def:physics:phyx-mini-0969:phyxminiproblems-problemphyxmini0969-angularspeedinradianspersecond, def:physics:phyx-mini-0969:phyxminiproblems-problemphyxmini0969-vacuumpermeabilityinnewtonsperamperesquared, def:physics:phyx-mini-0969:phyxminiproblems-problemphyxmini0969-rotatingchargedcylindricalshellsetup}
  Positivity and nondegeneracy conditions implicit in a physical shell
\end{definition}
\begin{proof}
  This is the declared model definition of Has Physical Rotating Charged Cylinder Parameters.
\end{proof}
\begin{definition}[Uses Coherent SIVacuum Permeability]
  \label{def:physics:phyx-mini-0969:phyxminiproblems-problemphyxmini0969-usescoherentsivacuumpermeability}
  \lean{PhyXMiniProblems.ProblemPhyXMini0969.UsesCoherentSIVacuumPermeability}
  \uses{def:physics:phyx-mini-0969:phyxminiproblems-problemphyxmini0969-vacuumpermeabilityinnewtonsperamperesquared, def:physics:phyx-mini-0969:phyxminiproblems-problemphyxmini0969-rotatingchargedcylindricalshellsetup}
  Physlib's EMSystem.μ₀ is a coherent-coordinate scalar. This assumption calibrates it to the separate dimensionful vacuum-permeability quantity
\end{definition}
\begin{proof}
  This is the declared model definition of Uses Coherent SIVacuum Permeability.
\end{proof}
\begin{definition}[Satisfies Uniform Rotating Surface Charge Model]
  \label{def:physics:phyx-mini-0969:phyxminiproblems-problemphyxmini0969-satisfiesuniformrotatingsurfacechargemodel}
  \lean{PhyXMiniProblems.ProblemPhyXMini0969.SatisfiesUniformRotatingSurfaceChargeModel}
  \uses{def:physics:phyx-mini-0969:phyxminiproblems-problemphyxmini0969-lengthinmeters, def:physics:phyx-mini-0969:phyxminiproblems-problemphyxmini0969-chargeincoulombs, def:physics:phyx-mini-0969:phyxminiproblems-problemphyxmini0969-angularspeedinradianspersecond, def:physics:phyx-mini-0969:phyxminiproblems-problemphyxmini0969-currentinamperes, def:physics:phyx-mini-0969:phyxminiproblems-problemphyxmini0969-surfacechargedensityincoulombspersquaremeter, def:physics:phyx-mini-0969:phyxminiproblems-problemphyxmini0969-linearchargedensityincoulombspermeter, def:physics:phyx-mini-0969:phyxminiproblems-problemphyxmini0969-ringcurrentperlengthinamperespermeter, def:physics:phyx-mini-0969:phyxminiproblems-problemphyxmini0969-rotatingchargedcylindricalshellsetup}
  Uniform-charge and rotation laws. The curved surface area is 2πRW, the charge per axial length is 2πRσ = Q/W, and a ring of charge dq rotating with angular speed ω carries current dI = ω dq/(2π). The final field is not mentioned here
\end{definition}
\begin{proof}
  This is the declared model definition of Satisfies Uniform Rotating Surface Charge Model.
\end{proof}
\begin{definition}[Satisfies Ring Biot Savart And Superposition]
  \label{def:physics:phyx-mini-0969:phyxminiproblems-problemphyxmini0969-satisfiesringbiotsavartandsuperposition}
  \lean{PhyXMiniProblems.ProblemPhyXMini0969.SatisfiesRingBiotSavartAndSuperposition}
  \uses{def:physics:phyx-mini-0969:phyxminiproblems-problemphyxmini0969-lengthinmeters, def:physics:phyx-mini-0969:phyxminiproblems-problemphyxmini0969-ringcurrentperlengthinamperespermeter, def:physics:phyx-mini-0969:phyxminiproblems-problemphyxmini0969-vacuumpermeabilityinnewtonsperamperesquared, def:physics:phyx-mini-0969:phyxminiproblems-problemphyxmini0969-magneticfluxdensityinteslas, def:physics:phyx-mini-0969:phyxminiproblems-problemphyxmini0969-axisvector, def:physics:phyx-mini-0969:phyxminiproblems-problemphyxmini0969-rotatingchargedcylindricalshellsetup}
  For a circular ring at axial coordinate x, Biot--Savart gives the axial field kernel dB/dx = μ₀ (dI/dx) R² / (2 (R² + x²)\textasciicircum{}(3/2)). The center field is the superposition of all rings from -W/2 to W/2. These hypotheses retain the unevaluated integrals and hence do not assume the closed form requested by the problem
\end{definition}
\begin{proof}
  This is the declared model definition of Satisfies Ring Biot Savart And Superposition.
\end{proof}
\begin{lemma}[rotating Cylinder Ring Kernel Integral]
  \label{lem:physics:phyx-mini-0969:phyxminiproblems-problemphyxmini0969-rotatingcylinderringkernelintegral}
  \lean{PhyXMiniProblems.ProblemPhyXMini0969.rotatingCylinderRingKernelIntegral}
  The calculus identity used after substituting the uniform charge and rotation relations. It is a conclusion to be proved, rather than a physical premise
\end{lemma}
\begin{proof}
  The conclusion follows from the governing relations and figure readouts named in the statement of rotating Cylinder Ring Kernel Integral.
\end{proof}
\begin{definition}[Answer Choice]
  \label{def:physics:phyx-mini-0969:phyxminiproblems-problemphyxmini0969-answerchoice}
  \lean{PhyXMiniProblems.ProblemPhyXMini0969.AnswerChoice}
  The four answer labels printed with the source problem
\end{definition}
\begin{proof}
  This is the declared model definition of Answer Choice.
\end{proof}
\begin{definition}[Displayed Magnetic Field Expression]
  \label{def:physics:phyx-mini-0969:phyxminiproblems-problemphyxmini0969-displayedmagneticfieldexpression}
  \lean{PhyXMiniProblems.ProblemPhyXMini0969.DisplayedMagneticFieldExpression}
  Choice C is printed as a scalar, while A, B, and D carry an explicit i-hat. This sum type preserves that distinction rather than silently turning every choice into a vector
\end{definition}
\begin{proof}
  This is the declared model definition of Displayed Magnetic Field Expression.
\end{proof}
\begin{definition}[displayed Expression]
  \label{def:physics:phyx-mini-0969:phyxminiproblems-problemphyxmini0969-answerchoice-displayedexpression}
  \lean{PhyXMiniProblems.ProblemPhyXMini0969.AnswerChoice.displayedExpression}
  \uses{def:physics:phyx-mini-0969:phyxminiproblems-problemphyxmini0969-lengthinmeters, def:physics:phyx-mini-0969:phyxminiproblems-problemphyxmini0969-chargeincoulombs, def:physics:phyx-mini-0969:phyxminiproblems-problemphyxmini0969-angularspeedinradianspersecond, def:physics:phyx-mini-0969:phyxminiproblems-problemphyxmini0969-rotatingchargedcylindricalshellsetup, def:physics:phyx-mini-0969:phyxminiproblems-problemphyxmini0969-answerchoice, def:physics:phyx-mini-0969:phyxminiproblems-problemphyxmini0969-displayedmagneticfieldexpression}
  The magnetic-field expression printed beside each answer label
\end{definition}
\begin{proof}
  This is the declared model definition of displayed Expression.
\end{proof}
\begin{definition}[recorded Dataset Answer]
  \label{def:physics:phyx-mini-0969:phyxminiproblems-problemphyxmini0969-recordeddatasetanswer}
  \lean{PhyXMiniProblems.ProblemPhyXMini0969.recordedDatasetAnswer}
  \uses{def:physics:phyx-mini-0969:phyxminiproblems-problemphyxmini0969-answerchoice}
  The answer label recorded by the source dataset
\end{definition}
\begin{proof}
  This is the declared model definition of recorded Dataset Answer.
\end{proof}
\begin{definition}[Matches Modeled Center Field]
  \label{def:physics:phyx-mini-0969:phyxminiproblems-problemphyxmini0969-matchesmodeledcenterfield}
  \lean{PhyXMiniProblems.ProblemPhyXMini0969.MatchesModeledCenterField}
  \uses{def:physics:phyx-mini-0969:phyxminiproblems-problemphyxmini0969-magneticfluxdensityinteslas, def:physics:phyx-mini-0969:phyxminiproblems-problemphyxmini0969-axisvector, def:physics:phyx-mini-0969:phyxminiproblems-problemphyxmini0969-rotatingchargedcylindricalshellsetup, def:physics:phyx-mini-0969:phyxminiproblems-problemphyxmini0969-answerchoice}
  A displayed expression agrees with the independent modeled center field
\end{definition}
\begin{proof}
  This is the declared model definition of Matches Modeled Center Field.
\end{proof}
% --- Archon declaration topology end ---
% --- Archon physics formalization source end ---
